\section*{Capítulo \ref{ch:b1:optical-mineralogy} --- Mineralogia óptica}

\begin{solution}{exo:b1:optical-mineralogy:1}
Quartzo: $\qty{30000}{nm}\times0.009 = \qty{270}{nm}$, um branco de primeira \omterm{def:b1:optical-mineralogy:interference-colour}{ordem}.
Calcita: $\qty{30000}{nm}\times0.172 = \qty{5160}{nm}$, bem além da carta:
um branco pálido de \omterm{def:b1:optical-mineralogy:interference-colour}{ordem} alta, levemente rosado ou verde.
\end{solution}

\begin{solution}{exo:b1:optical-mineralogy:2}
$I/I_{\max} = \sin^2 2\theta$: $0$ em \ang{0}; $\sin^2\ang{45} = 0.5$ em
\ang{22.5}; $1$ em \ang{45}; $0.5$ em \ang{67.5}; $0$ em \ang{90}.
\end{solution}

\begin{solution}{exo:b1:optical-mineralogy:3}
A \omterm{def:b1:rock-forming-minerals:groups}{granada} é cúbica, logo isotrópica: a luz sai dela vibrando como entrou,
ao longo do polarizador, e o analisador a detém; $\Delta = 0$ em toda posição. Um grão de quartzo cortado perpendicularmente ao seu \omterm{def:b1:optical-mineralogy:optic-axis}{eixo óptico} também tem
$n_1 = n_2$ para a luz que viaja ao longo desse eixo, logo $\Delta = 0$: ele também permanece escuro. Em \omterm{def:b1:field-lab-methods-1:microscope}{luz plano-polarizada} a \omterm{def:b1:rock-forming-minerals:groups}{granada} tem \omterm{def:b1:optical-mineralogy:relief}{relevo} muito alto ($1.83$
contra $1.54$), o quartzo quase nenhum.
\end{solution}

\begin{solution}{exo:b1:optical-mineralogy:4}
$|n - 1.54|$: almandina $0.29$, forsterita $0.11$, ortoclásio $0.02$, quartzo
$0.01$; \omterm{def:b1:optical-mineralogy:relief}{relevo} nessa \omterm{def:b1:optical-mineralogy:interference-colour}{ordem}, do mais alto ao mais baixo. O ortoclásio tem índice menor: \omterm{def:b1:optical-mineralogy:relief}{relevo} negativo, e a sua \omterm{def:b1:optical-mineralogy:relief}{linha de Becke} sai do grão, para a montagem, quando o foco é elevado.
\end{solution}

\begin{solution}{exo:b1:optical-mineralogy:5}
$\delta = 650/30000 = 0.022$. Um \omterm{def:b1:rock-forming-minerals:mineral}{mineral} verde \omterm{def:b1:optical-mineralogy:pleochroism}{pleocroico} com clivagens a \ang{56} e \ang{124} e \omterm{def:b1:optical-mineralogy:birefringence}{birrefringência} próxima de $0.020$: hornblenda.
\end{solution}

\begin{solution}{exo:b1:optical-mineralogy:6}
$\tan\ang{30} = 0.5774$. Quartzo: $(0.013/1.54)\times0.5774 = \qty{0.00487}{rad}
= \ang{0.28}$. Forsterita: $(0.111/1.54)\times0.5774 = \qty{0.0416}{rad} =
\ang{2.38}$. Razão $0.111/0.013 = 8.5$. Exatamente: $\sin r = 1.54\sin\ang{30}/1.651
= 0.4664$, $r = \ang{27.80}$, um desvio de \ang{2.20}: a fórmula de primeira \omterm{def:b1:optical-mineralogy:interference-colour}{ordem} superestima em \qty{8}{\percent} para uma diferença de índice tão grande.
\end{solution}

\begin{solution}{exo:b1:optical-mineralogy:7}
$270 + 540 = \qty{810}{nm}$, um verde de segunda \omterm{def:b1:optical-mineralogy:interference-colour}{ordem}, quando as direções lentas do grão e da lâmina são paralelas (os \omterm{def:b1:optical-mineralogy:retardation}{retardos} se somam); $|270 - 540| =
\qty{270}{nm}$, branco de primeira \omterm{def:b1:optical-mineralogy:interference-colour}{ordem} de novo, quando a direção lenta do grão é paralela à direção rápida da lâmina.
\end{solution}

\begin{solution}{exo:b1:optical-mineralogy:8}
Na espessura total, $\Delta = \qty{30000}{nm}\times0.052 = \qty{1560}{nm}$. Da borda ($\Delta = 0$) para dentro, o \omterm{def:b1:optical-mineralogy:retardation}{retardo} passa por 542 e
\qty{1074}{nm}, mas não por \qty{1605}{nm}: duas franjas vermelho-violeta.
\end{solution}

\begin{solution}{exo:b1:optical-mineralogy:9}
O \omterm{def:b1:optical-mineralogy:extinction}{ângulo de extinção} medido em um grão depende de como ele é cortado: só seções de uma orientação particular mostram o maior ângulo, e outras mostram menos, até zero. O maior valor, cerca de \ang{37}, é o que deve ser comparado: ele coincide com o diopsídio (cerca de \ang{38}), e não com a hornblenda (no máximo \ang{19}).
\end{solution}

\begin{solution}{exo:b1:optical-mineralogy:10}
Com o analisador ao longo de $x$, as componentes em $x$ das duas ondas são
$E_0\cos\theta\cdot\cos\theta$ e $-E_0\sin\theta\,e^{-i\varphi}\cdot(-\sin\theta)$,
logo $A = E_0(\cos^2\theta + \sin^2\theta\,e^{-i\varphi})$ e
\[
  |A|^2 = E_0^2\bigl(\cos^4\theta + \sin^4\theta + 2\sin^2\theta\cos^2\theta\cos\varphi\bigr)
  = E_0^2\bigl(1 - 2\sin^2\theta\cos^2\theta\,(1 - \cos\varphi)\bigr),
\]
usando $\cos^4\theta + \sin^4\theta = 1 - 2\sin^2\theta\cos^2\theta$. Como
$2\sin^2\theta\cos^2\theta = \frac12\sin^2 2\theta$ e $1 - \cos\varphi =
2\sin^2(\varphi/2)$, $I_\parallel = I_0(1 - \sin^2 2\theta\sin^2(\pi\Delta/\lambda))$.
Em cada comprimento de onda, $I_\parallel + I_\perp = I_0$: as cores entre polarizadores paralelos são os complementos das que aparecem entre \omterm{def:b1:field-lab-methods-1:microscope}{polarizadores cruzados}. Um grão de \qty{550}{nm}, magenta entre \omterm{def:b1:field-lab-methods-1:microscope}{polarizadores cruzados}, é verde-amarelado pálido entre polarizadores paralelos.
\end{solution}

\begin{solution}{exo:b1:optical-mineralogy:11}
(a) $\Delta/\lambda$ vai de $5000/700 = 7.14$ a $5000/400 = 12.5$; passa pelos inteiros 8 a 12: cinco zeros. (b) Os zeros estão em $\lambda_k = \Delta/k$; dois sucessivos distam $\Delta/k - \Delta/(k + 1) =
\Delta/(k(k + 1)) \approx \lambda^2/\Delta$, isto é, $550^2/5000 =
\qty{60.5}{nm}$ perto de \qty{550}{nm}. (c) Cinco faixas escuras e estreitas espalhadas pelo espectro retiram um pouco de cada cor, e o restante passa, em média, com metade da intensidade: o olho vê um tom pálido, quase neutro.
\end{solution}

\begin{solution}{exo:b1:optical-mineralogy:12}
A média de $\sin^2 2\theta$ em uma volta é $\frac1{2\pi}\int_0^{2\pi}\sin^2 2\theta\,
\mathrm d\theta = \frac12$, logo a intensidade média é
$\frac12 I_0\sin^2(\pi\Delta/\lambda)$. Ela calcula $\delta = \Delta/30000$
em vez de $\Delta/25000$ (em nm): todos os seus valores são $25/30 = 0.83$ dos
verdadeiros, baixos em \qty{17}{\percent}.
\end{solution}

\begin{solution}{pb:b1:optical-mineralogy:1}
\textbf{1.} B: as bordas escuras significam \omterm{def:b1:optical-mineralogy:relief}{relevo} alto; as fronteiras quase invisíveis de A significam \omterm{def:b1:optical-mineralogy:relief}{relevo} baixo.

\textbf{2.} A: $(1.544 + 1.553)/2 = 1.5485$, $|n - n_m| = 0.0085$. B:
$(1.664 + 1.672 + 1.694)/3 = 1.6767$, $|n - n_m| = 0.1367$.

\textbf{3.} Os desvios são proporcionais a $|n - n_m|$: $0.1367/0.0085 =
16.1$.

\textbf{4.} B tem índice mais alto que a montagem: \omterm{def:b1:optical-mineralogy:relief}{relevo} positivo.

\textbf{5.} \omterm{def:b1:optical-mineralogy:pleochroism}{Pleocroísmo}; biotita.

\textbf{6.} Duas clivagens a \ang{56} e \ang{124}: um \omterm{def:b1:rock-forming-minerals:groups}{anfibólio}.

\textbf{7.} $\qty{30000}{nm}\times0.009 = \qty{270}{nm}$: branco de primeira \omterm{def:b1:optical-mineralogy:interference-colour}{ordem}.

\textbf{8.} Cada grão é cortado na sua própria orientação e mostra um valor de
$|n_1 - n_2|$ entre 0 e $\delta$; só os grãos cortados paralelamente ao \omterm{def:b1:optical-mineralogy:optic-axis}{eixo óptico} mostram o próprio $\delta$. A cor mais alta entre muitos grãos é a que deve ser usada.

\textbf{9.} $d = 342/0.009 = \qty{38000}{nm} = \qty{38}{\micro m}$.

\textbf{10.} $(38 - 30)/30 = \qty{26.7}{\percent}$.

\textbf{11.} $\sin^2(\pi\times342/550) = 0.861$; em \qty{30}{\micro m},
$\sin^2(\pi\times270/550) = 0.999$.

\textbf{12.} $\qty{38000}{nm}\times0.008 = \qty{304}{nm}$ até
$\qty{38000}{nm}\times0.013 = \qty{494}{nm}$: de branco a vermelho-laranja, tudo abaixo
do fim da primeira \omterm{def:b1:optical-mineralogy:interference-colour}{ordem} em \qty{542}{nm}.

\textbf{13.} Entre 1074 e \qty{1605}{nm}: terceira \omterm{def:b1:optical-mineralogy:interference-colour}{ordem}.

\textbf{14.} $\delta = 1140/38000 = 0.030$.

\textbf{15.} $1140/30000 = 0.038$, o que sugeriria uma \omterm{def:b1:rock-forming-minerals:groups}{olivina}. Uma \omterm{def:b1:rock-forming-minerals:groups}{olivina} magnesiana e o quartzo reagem para dar \omterm{def:b1:rock-forming-minerals:groups}{piroxênio},
\ce{Mg2SiO4 + SiO2 -> 2 MgSiO3}; uma rocha que cristalizou com quartzo não deve manter ambos.

\textbf{16.} $0.030$ é a \omterm{def:b1:optical-mineralogy:birefringence}{birrefringência} do diopsídio; seu \omterm{def:b1:optical-mineralogy:relief}{relevo} alto, sua falta de cor e suas duas clivagens perto de \ang{87} concordam: B é um clinopiroxênio, diopsídio.

\textbf{17.} $\qty{30000}{nm}\times0.030 = \qty{900}{nm}$: amarelo de segunda \omterm{def:b1:optical-mineralogy:interference-colour}{ordem}.

\textbf{18.} Por \qty{8}{\micro m}; o quartzo então mostra no máximo \qty{270}{nm}, branco de primeira \omterm{def:b1:optical-mineralogy:interference-colour}{ordem}.

\textbf{19.} \omterm{def:b1:optical-mineralogy:twin}{Maclagem polissintética}: D é \omterm{def:b1:rock-forming-minerals:groups}{plagioclásio}.

\textbf{20.} A composição do \omterm{def:b1:rock-forming-minerals:groups}{plagioclásio} (o seu teor de anortita).

\textbf{21.} Um \omterm{def:b1:rock-forming-minerals:groups}{anfibólio} com \omterm{def:b1:optical-mineralogy:extinction}{ângulo de extinção} de \ang{17}: hornblenda
(\ang{15}--\ang{19}), não um \omterm{def:b1:rock-forming-minerals:groups}{piroxênio} (\ang{38} para o diopsídio).

\textbf{22.} $\qty{38000}{nm}\times0.020 = \qty{760}{nm}$: um verde de segunda \omterm{def:b1:optical-mineralogy:interference-colour}{ordem}.

\textbf{23.} Ele é cortado perpendicularmente a um \omterm{def:b1:optical-mineralogy:optic-axis}{eixo óptico}: os seus dois índices são iguais para a luz que viaja ao longo do eixo, $\Delta = 0$, e ele permanece escuro como um grão isotrópico.

\textbf{24.} Quartzo A, diopsídio B, biotita C, \omterm{def:b1:rock-forming-minerals:groups}{plagioclásio} D, hornblenda E: uma rocha ígnea intermediária com algum quartzo, como um quartzo diorito.

\textbf{25.} A seção tem \textbf{\qty[text-series-to-math=true]{38}{\micro m}} de espessura.
\end{solution}
